% GENERATED FILE. Edit canonical lessons, then run npm run generate:books.
% canonical-source-hash: fnv1a64:d68db1dcc0c73b62
% canonical-lessons: JA-C486-buhin, JA-C486-kanagu, JA-C486-sabiru, JA-C486-hameru, JA-C486-hagasu

\chapter{Parts and Fittings}
\label{ch:ja-parts-and-fittings}
\hlchaptermodality{486}

\begin{chapteropening}
\textbf{By the end of this chapter:} \emph{I can say what is broken and what needs replacing.}

Complete the last lesson of chapter 486: I can say what is broken and what needs replacing.
\end{chapteropening}

\section[\texorpdfstring{\ja{ぶひん} (buhin)}{ぶひん (buhin)}]{\ja{ぶひん} (buhin) — a part (of a machine)}

\label{lesson:JA-C486-buhin}



\begin{quote}
Before the new one: say the Japanese for plain, then the Japanese for showy.
\end{quote}

\subsection*{You'll want to know: \ja{ぶひん}}
\textbf{\ja{ぶひん}} — \emph{buhin} — \textquotedblleft{}a part (of a machine)\textquotedblright{}.

\begin{quote}
\textbf{\ja{ぶひんが} \ja{こわれたら、とりかえます。}} — \emph{buhin ga kowaretara, torikaemasu} — \textquotedblleft{}If a part breaks, I replace it.\textquotedblright{}
\end{quote}

\begin{cousinweb}
\textbf{\ja{ぶひん}} (\emph{buhin}) is Sino-Japanese (kango): two Chinese characters read with their on'yomi, readings borrowed from Chinese. They are \emph{bu}, \textquotedblleft{}section\textquotedblright{} (Middle Chinese \emph{buX}) and \emph{hin}, \textquotedblleft{}goods\textquotedblright{} (\emph{phimX}) (KANJIDIC2; Middle Chinese in Baxter's transcription).
\end{cousinweb}

\begin{grammarlens}[title={What you've built}]
One more word for describing what you lived and felt.
\end{grammarlens}

\subsection*{Your turn}
\begin{itemize}
\raggedright
  \item \emph{Say it:} \emph{buhin}
  \item \emph{Say it:} \emph{buhin}, once more
  \item \emph{Say it:} \emph{buhin ga kowaretara, torikaemasu} — If a part breaks, I replace it.
  \item \emph{Recall:} say the Japanese for plain, then the Japanese for showy, then say \emph{buhin} again
\end{itemize}

\begin{tcolorbox}[breakable,colback=teal!4,colframe=teal!35!black,title={Before you move on}]
What does \ja{ぶひん} mean? (\textquotedblleft{}A part (of a machine)\textquotedblright{}.) Say it once more.
\end{tcolorbox}
\section[\texorpdfstring{\ja{かなぐ} (kanagu)}{かなぐ (kanagu)}]{\ja{かなぐ} (kanagu) — a metal fitting}

\label{lesson:JA-C486-kanagu}



\begin{quote}
Before the new one: say the Japanese for showy, then the Japanese for a part (of a machine).
\end{quote}

\subsection*{You'll want to know: \ja{かなぐ}}
\textbf{\ja{かなぐ}} — \emph{kanagu} — \textquotedblleft{}a metal fitting\textquotedblright{}.

\begin{quote}
\textbf{\ja{かばんの} \ja{かなぐが} \ja{こわれました。}} — \emph{kaban no kanagu ga kowaremashita} — \textquotedblleft{}The metal fitting on my bag broke.\textquotedblright{}
\end{quote}

\begin{cousinweb}
\textbf{\ja{かなぐ}} (\emph{kanagu}) mixes the two layers, native and Sino-Japanese: native \emph{kana}, \textquotedblleft{}gold\textquotedblright{} and Sino-Japanese \emph{gu}, \textquotedblleft{}tool\textquotedblright{} (Middle Chinese \emph{gjuH}), borrowed from Chinese (JMdict, KANJIDIC2; Middle Chinese in Baxter's transcription).
\end{cousinweb}

\begin{grammarlens}[title={What you've built}]
One more word for describing what you lived and felt.
\end{grammarlens}

\subsection*{Your turn}
\begin{itemize}
\raggedright
  \item \emph{Say it:} \emph{kanagu}
  \item \emph{Say it:} \emph{kanagu}, once more
  \item \emph{Say it:} \emph{kaban no kanagu ga kowaremashita} — The metal fitting on my bag broke.
  \item \emph{Recall:} say the Japanese for showy, then the Japanese for a part (of a machine), then say \emph{kanagu} again
\end{itemize}

\begin{tcolorbox}[breakable,colback=teal!4,colframe=teal!35!black,title={Before you move on}]
What does \ja{かなぐ} mean? (\textquotedblleft{}A metal fitting\textquotedblright{}.) Say it once more.
\end{tcolorbox}
\section[\texorpdfstring{\ja{さびる} (sabiru)}{さびる (sabiru)}]{\ja{さびる} (sabiru) — to rust}

\label{lesson:JA-C486-sabiru}



\begin{quote}
Before the new one: say the Japanese for a part (of a machine), then the Japanese for a metal fitting.
\end{quote}

\subsection*{You'll want to know: \ja{さびる}}
\textbf{\ja{さびる}} — \emph{sabiru} — \textquotedblleft{}to rust\textquotedblright{}.

\begin{quote}
\textbf{\ja{じてんしゃが} \ja{あめで} \ja{さびました。}} — \emph{jitensha ga ame de sabimashita} — \textquotedblleft{}My bicycle rusted in the rain.\textquotedblright{}
\end{quote}

Polite \textbf{\ja{さびます}}, \emph{sabimasu}; an ichidan verb. Past (polite): \textbf{\ja{じてんしゃが} \ja{さびました。}}, \emph{jitensha ga sabimashita}.

\begin{cousinweb}
\textbf{\ja{さびる}} (\emph{sabiru}) is a native Japanese word (wago), inherited rather than borrowed from Chinese; its kanji was attached later as a kun'yomi, and its older origin is not settled (JMdict, KANJIDIC2).
\end{cousinweb}

\begin{grammarlens}[title={What you've built}]
One more word for describing what you lived and felt.
\end{grammarlens}

\subsection*{Your turn}
\begin{itemize}
\raggedright
  \item \emph{Say it:} \emph{sabiru}
  \item \emph{Say it:} \emph{sabiru}, once more
  \item \emph{Say it:} \emph{jitensha ga ame de sabimashita} — My bicycle rusted in the rain.
  \item \emph{Recall:} say the Japanese for a part (of a machine), then the Japanese for a metal fitting, then say \emph{sabiru} again
\end{itemize}

\begin{tcolorbox}[breakable,colback=teal!4,colframe=teal!35!black,title={Before you move on}]
What does \ja{さびる} mean? (\textquotedblleft{}To rust\textquotedblright{}.) Say it once more.
\end{tcolorbox}
\section[\texorpdfstring{\ja{はめる} (hameru)}{はめる (hameru)}]{\ja{はめる} (hameru) — to fit in; to put on (a ring)}

\label{lesson:JA-C486-hameru}



\begin{quote}
Before the new one: say the Japanese for a metal fitting, then the Japanese for to rust.
\end{quote}

\subsection*{You'll want to know: \ja{はめる}}
\textbf{\ja{はめる}} — \emph{hameru} — \textquotedblleft{}to fit in; to put on (a ring)\textquotedblright{}.

\begin{quote}
\textbf{\ja{ゆびわを} \ja{ゆびに} \ja{はめます。}} — \emph{yubiwa o yubi ni hamemasu} — \textquotedblleft{}I put the ring on my finger.\textquotedblright{}
\end{quote}

Polite \textbf{\ja{はめます}}, \emph{hamemasu}; an ichidan verb. Past (polite): \textbf{\ja{あねが} \ja{ゆびわを} \ja{はめました。}}, \emph{ane ga yubiwa o hamemashita}.

\begin{cousinweb}
\textbf{\ja{はめる}} (\emph{hameru}) is a native Japanese word (wago), inherited rather than borrowed from Chinese, and usually written in kana; its older origin is not settled (JMdict).
\end{cousinweb}

\begin{grammarlens}[title={What you've built}]
One more word for describing what you lived and felt.
\end{grammarlens}

\subsection*{Your turn}
\begin{itemize}
\raggedright
  \item \emph{Say it:} \emph{hameru}
  \item \emph{Say it:} \emph{hameru}, once more
  \item \emph{Say it:} \emph{yubiwa o yubi ni hamemasu} — I put the ring on my finger.
  \item \emph{Recall:} say the Japanese for a metal fitting, then the Japanese for to rust, then say \emph{hameru} again
\end{itemize}

\begin{tcolorbox}[breakable,colback=teal!4,colframe=teal!35!black,title={Before you move on}]
What does \ja{はめる} mean? (\textquotedblleft{}To fit in; to put on (a ring)\textquotedblright{}.) Say it once more.
\end{tcolorbox}
\section[\texorpdfstring{\ja{はがす} (hagasu)}{はがす (hagasu)}]{\ja{はがす} (hagasu) — to peel off; to tear off}

\label{lesson:JA-C486-hagasu}



\begin{quote}
Before the new one: say the Japanese for to rust, then the Japanese for to fit in.
\end{quote}

\subsection*{You'll want to know: \ja{はがす}}
\textbf{\ja{はがす}} — \emph{hagasu} — \textquotedblleft{}to peel off; to tear off\textquotedblright{}.

\begin{quote}
\textbf{\ja{ふうとうの} \ja{きってを} \ja{はがします。}} — \emph{fūtō no kitte o hagashimasu} — \textquotedblleft{}I peel the stamp off the envelope.\textquotedblright{}
\end{quote}

Polite \textbf{\ja{はがします}}, \emph{hagashimasu}; a godan verb. Past (polite): \textbf{\ja{こどもが} \ja{きってを} \ja{はがしました。}}, \emph{kodomo ga kitte o hagashimashita}.

\begin{cousinweb}
\textbf{\ja{はがす}} (\emph{hagasu}) is a native Japanese word (wago), inherited rather than borrowed from Chinese, and usually written in kana; its older origin is not settled (JMdict).
\end{cousinweb}

\begin{grammarlens}[title={What you've built}]
One more word for describing what you lived and felt.
\end{grammarlens}

\subsection*{Your turn}
\begin{itemize}
\raggedright
  \item \emph{Say it:} \emph{hagasu}
  \item \emph{Say it:} \emph{hagasu}, once more
  \item \emph{Say it:} \emph{fūtō no kitte o hagashimasu} — I peel the stamp off the envelope.
  \item \emph{Recall:} say the Japanese for to rust, then the Japanese for to fit in, then say \emph{hagasu} again
\end{itemize}

\begin{tcolorbox}[breakable,colback=teal!4,colframe=teal!35!black,title={Before you move on}]
What does \ja{はがす} mean? (\textquotedblleft{}To peel off; to tear off\textquotedblright{}.) Say it once more.
\end{tcolorbox}
